\section{The cubic coefficient and the exact scope of the common integer bound}
\label{sec:cubic-propgk-size}

This section reads the actual source interface behind the explanatory
name \texttt{PropGK}.  All Lean paths below are relative to
\path{src/Development/AutomorphicArithmetic/AutomorphicVarieties/}
at commit \texttt{7cbf7a0c32fca1badf05931879a37c3f8a2b03b4}.
Reading the source proofs is not a claim to have executed Lean or to
have independently verified every imported theorem.

\paragraph{The source statement and the point of generality.}
\path{UnitaryThreeGroups/KatoUse.lean}, definitions
\texttt{PropGK}, \texttt{CorGK}, and \texttt{OneChar}, quantifies
over \emph{Dirichlet} characters of $(\mathbb Z/N)^\times$.
The imaginary quadratic field $K$ is an auxiliary field for the
unitary argument.  More precisely,
\path{HeightAndVanishing/DirichletPrimeToLPartMulLPowerPart.lean},
\texttt{Vlen}, \texttt{Vn}, and \texttt{orb}, defines
\[
 \begin{split}
 V_{\ell,e,\psi}
   &=\operatorname{length}_{O}
       \bigl(\operatorname{ShaU}_{\ell N}
        ((\operatorname{Frac}O/O)(\psi^{-1})(-e))\bigr),
       \quad\hbox{over }G_{\mathbb Q},\\
 V^{\rm n}_{\ell,e,\psi}&=V_{\ell,e,\psi}/e_O,
 \qquad
 \operatorname{orb}(H)=\{\psi:\ker\psi=H\}.
 \end{split}                                                    \tag{GK.1}
\]
Here $e_O=\operatorname{length}_O(O/\ell O)$, and the precise
local-condition family is the source's \texttt{ShaU}; it must not
be silently replaced by another Selmer condition.  The definitions
of \texttt{Chr.rho}, \texttt{VS}, \texttt{V'}, and \texttt{nrm}
are in \path{CharacterLevelBounds/ShaFamMono.lean}.

The entire proof of
\path{ModularityArgument/ThresholdArguments/DenseThresholdCases/FPropGKU.lean},
\texttt{propGK\_of\_numberU3U\_of\_oneCharU}, has been read.
It separates into two inputs.  The one-character input extracts
finitely many automorphic trace differences from the cohomology
length.  The other input, \texttt{NumberU3U\_s}, bounds their product
by one nonzero integer, chosen before $\ell$.
Unlike \texttt{PropGK}, this latter input permits any finite,
Galois-stable family of roots of unity.  Its exact definition is in
\path{ModularityArgument/ThresholdArguments/DenseThresholdCases/NumberU3US.lean}.
The complete assembly proof in \texttt{KeepNumU.lean},
\texttt{numberU3U\_of\_adm2A\_imp\_adm2B}, combines fixed-$K$
cuspidal trace data and rational rank-two trace data.  If their
counting constants are $(c_1,A_1)$ and $(c_2,A_2)$, it gives
\[
 c=5\bigl(c_1+c_2|D_K|^{A_2}\bigr),\qquad A=A_1+A_2+1.
                                                               \tag{GK.2}
\]
Those counting and automorphy inputs remain source inputs here.

For clarity, denote its unchanged admissible class by
$\mathcal A_{\ell,e,M}$: it is the union of
\texttt{Adm3U\_s} at level $M$ and \texttt{Adm2A} at level
$M|D_K|$, with exactly the source's attachment, local-level and
safety conditions.  Put
\[
 E_{q,e}(z)=q^{-e}z+1+q^ez^{-1},\qquad v_\ell(\ell)=1.
                                                               \tag{GK.3}
\]
The input \texttt{NumberU3U\_s} says that, for fixed quadratic
$K$, there are $c,A$ such that for $e\ge1$, $M\ge1$,
$p\ge7$ prime with $p\nmid M D_K$, a place $w\mid p$, and
a finite Galois-stable root-of-unity family $(z_o)_{o\in\mathcal O}$,
there is $J_p\in\mathbb Z\setminus\{0\}$ with
\[
 \log|J_p|\le c|\mathcal O|[M(e+1)]^A\log p.                    \tag{GK.4}
\]
For every $\ell\ne2,p$, every choice of the source coefficient
embedding, and every finite family of pairwise nonisomorphic
members of $\mathcal A_{\ell,e,M}$ for each $o$, all trace
differences at $\operatorname{Frob}_w^{-1}$ from
$E_{Nw,e}(z_o)$ are nonzero and
\[
 \sum_o\sum_{\rho}v_\ell
  \bigl(\operatorname{tr}\rho(\operatorname{Frob}_w^{-1})
                       -E_{Nw,e}(z_o)\bigr)\le v_\ell(J_p).
                                                               \tag{GK.5}
\]
Complex algebraic quantities in this formula are transported to
$\overline{\mathbb Q}_\ell$ by that embedding.
Equation (GK.5) is the source's norm inequality after taking
$-\log_\ell$, so its direction and its normalization are fixed.
The complete proof of
\path{BuildingsResidueChars/FrobeniusTraceNormNonzeroLogBound.lean}
was read: its integer is the Galois-invariant product of
$q^e(t_s-E_{q,e}(z_o))$ over the full finite trace data.
Its literal \texttt{Af.TrData} input uses $|t_s|\le3$ and integrality of $q^et_s$; the separate \texttt{Win.TrData} has the weaker bound $\sqrt q+2$. AC-G112 and AC-G112-twist explicitly supply the stronger \texttt{Af} fields. In the product argument,
omitting factors can only increase the $\ell$-adic norm when
$\ell\nmid q$.

\begin{proposition}[AC-K29: controlled inert primes and a conditional cubic size bound]
\label{AC-K29}
Let $K=\mathbb Q(\omega)$ and let
$\chi:G_K\to\mu_3$ be the character of
$K(\sqrt[3]{269})/K$.  Let $N\mid M$, with $269\mid M$.
Fix $e\ge1$ and a finite family
$\mathcal O=\{(\psi_o,j_o)\}$, where $\psi_o$ is a Dirichlet
character modulo $N$ and $j_o\in\{0,1,2\}$.
Assume the family is stable, with multiplicities, under automorphisms
of its algebraic coefficient field.  Define
\[
 R_o=(\psi_o|_{G_K}\chi^{j_o}\kappa^e)
       \oplus1\oplus
       (\psi_o^{-1}|_{G_K}\chi^{j_o}\kappa^{-e}),              \tag{GK.6}
\]
where $\kappa$ is the cyclotomic character.
There are distinct primes $p,p'$ such that
\[
 \begin{gathered}
 p\equiv p'\equiv2\pmod3,
 \qquad11\le p<30M,\quad11\le p'<(30M)^2,\\
 pp'\text{ is prime to }30M,
 \qquad\log p+\log p'\le3\log(30M),\qquad
 \chi(\operatorname{Frob}_w)=1\quad(w\mid p\text{ or }p').
 \end{gathered}                                                \tag{GK.7}
\]
The last assertion refers to the arithmetic Frobenius at the unique
place of $K$ over either inert prime.

Assume the fixed-$K$ source input \texttt{NumberU3U\_s}, as
specified in (GK.4)--(GK.5).  Let $\mathcal E$ be any set of
eligible coefficient primes, all different from $2$, for which
the following \emph{cubic extraction hypothesis} holds.
For each $\ell\in\mathcal E$, $o\in\mathcal O$, and either
auxiliary prime $r\in\{p,p'\}\setminus\{\ell\}$, use the actual
strict Selmer group of AC-K25 with coefficient
$(\operatorname{Frac}O_\ell/O_\ell)
 (\psi_o^{-1}|_{G_K}\chi^{-j_o})(-e)$.
It has finite normalized length
\[
 s_{\ell,o}=\frac{\operatorname{length}_{O_\ell}S_{\ell,o}}
                   {\operatorname{length}_{O_\ell}(O_\ell/\ell O_\ell)}
       =\frac{\log\#S_{\ell,o}}{[\operatorname{Frac}O_\ell:
                              \mathbb Q_\ell]\log\ell}\ge0.
                                                               \tag{GK.8a}
\]
There is a finite family $X_{\ell,o,r}$
of representations of $G_K$, and a comparison to members of
$\mathcal A_{\ell,e,M}$, which is injective on isomorphism
classes, preserves traces at $\operatorname{Frob}_{w_r}^{-1}$,
and satisfies
\[
 s_{\ell,o}\le2\sum_{\rho\in X_{\ell,o,r}}v_\ell
  \left(\operatorname{tr}\rho(\operatorname{Frob}_{w_r}^{-1})
             -\operatorname{tr}R_o(\operatorname{Frob}_{w_r}^{-1})\right).
                                                               \tag{GK.8}
\]
No isomorphism class is repeated within $X_{\ell,o,r}$.
Then
\[
 \sum_{\ell\in\mathcal E}\log\ell\sum_{o\in\mathcal O}s_{\ell,o}
   \le6c|\mathcal O|[M(e+1)]^A\log(30M).                     \tag{GK.9}
\]
In particular the left side has finite support.  This assertion
bounds only the eligible primes for which (GK.8) and its comparison
have been proved; it does not include excluded primes by convention.
\end{proposition}
\begin{proof}
The integer $30M-1$ is congruent to $2$ modulo $3$.
At least one of its prime factors is therefore $2$ modulo $3$.
Choose such a factor $p$.  It divides neither $M$ nor $2,3,5$,
so $p\ge11$ and $p<30M$.  Applying the same argument to
$30Mp-1$ gives $p'\ne p$, with all the asserted bounds and
coprimalities.  Their logarithmic bound follows immediately.

For $r=p$ or $p'$, the cube map on $\mathbb F_r^\times$ is a
bijection, since $\gcd(3,r-1)=1$.  Thus $X^3-269$ has a root
in $\mathbb F_r$.  Its other two roots lie in $\mathbb F_{r^2}$,
which contains $\mu_3$.  Since $r\nmid3\cdot269$, the extension
is unramified there and its Frobenius over $K$ is trivial.
This proves the last part of (GK.7).

Put $q=r^2$ and $z_{o,r}=\psi_o(q)^{-1}$.
Restriction of a Dirichlet character and of $\kappa$ to the
arithmetic Frobenius over $K$ evaluates them as $\psi_o(q)$ and
$q$, respectively.  Consequently (GK.6) gives the exact equality
\[
 \operatorname{tr}R_o(\operatorname{Frob}_{w_r}^{-1})
                        =E_{q,e}(z_{o,r}).                    \tag{GK.10}
\]
The family $z_{o,r}$ consists of roots of unity and is Galois
stable.  Applying (GK.4)--(GK.5) at $p$ and at $p'$ produces
integers $J_p,J_{p'}$ independent of $\ell$.  Use the
trace-preserving comparison in (GK.8) to obtain
\[
 \sum_os_{\ell,o}\le2v_\ell(J_r)
        \quad(\ell\in\mathcal E,\ \ell\ne r).                \tag{GK.11}
\]
The differences are nonzero by the same source input, so all
valuations used here are finite.
For any finite subset of $\mathcal E$, use $r=p$ at primes other
than $p$, and use $r=p'$ if $\ell=p$.  The factorization of a
nonzero rational integer yields
\[
 \sum_{\ell\in\mathcal E}\log\ell\sum_os_{\ell,o}
   \le2\log|J_p|+2\log|J_{p'}|.
\]
One may first read this as a bound for every finite subset.
Equation (GK.11) also makes every positive summand supported on
the prime divisors of $J_pJ_{p'}$, proving finite support directly.
Equations (GK.4) and (GK.7) now give (GK.9).
\end{proof}

\begin{corollary}[AC-K29-budget: the conductor and weight scale of this reduction]
\label{AC-K29-budget}
Suppose the extraction hypotheses of AC-K29 are established over
the actual family, with $|\mathcal O_x|\le D_x$ and one level
$M_x$ for every $1\le e<t_x$.  Then its eligible-prime sum is at most
\[
 6cD_x M_x^A t_x(t_x+1)^A\log(30M_x).                       \tag{GK.12}
\]
If $t_x\le x^{401}$, $D_x$ is bounded by a fixed power of $x$,
and
\[
 \log M_x=O_\varepsilon(x^{\varepsilon/5}(\log x)^2),          \tag{GK.13}
\]
this bound is $o(\exp(x^\varepsilon/2))$.
A fixed power of the conductor in $M_x$ is allowed in (GK.13).
The source constants $c,A$ depend on the fixed field $K$,
not on $x,e$ or the coefficient prime.
\end{corollary}
\begin{proof}
Sum (GK.9) and use
$\sum_{1\le e<t}(e+1)^A\le t(t+1)^A$.
The logarithm of (GK.12), after replacing a zero factor by one
if necessary, is at most
$A\log M_x+O_\varepsilon(\log x+\log\log(30M_x))=o(x^\varepsilon)$.
Subtracting $x^\varepsilon/2$ proves the claimed little-o estimate.
Replacing a conductor by any fixed power multiplies its logarithm
by a fixed constant and preserves this argument.
\end{proof}

\paragraph{Why this is a new extraction obligation.}
Complex conjugation gives $\chi^c=\chi^{-1}$.
For $j=1,2$, neither $\chi^j$ nor $\chi^j\kappa^e$ is the
restriction of a one-dimensional character of $G_{\mathbb Q}$:
such restrictions are fixed by conjugation, whereas equality here
would require $\chi^{2j}=1$.
The CM Hom realization of AC-G23 does not change this fact.
The corrected packet (GK.6) is conjugate self-dual, because
conjugate duality exchanges its two endpoints.  Its determinant,
however, is $\chi^{2j}$, and its middle character is exactly $1$.
It is therefore not the source's rational packet
$\Psi\oplus1\oplus\Psi^{-1}$ with determinant one.

The entire proof of \texttt{FOneCharU.lean},
\texttt{oneCharU\_of\_katoTwoW\_rig\_admI\_automorphy}, was read.
Its \texttt{Dat}, \texttt{exists\_dat}, \texttt{model\_upper},
\texttt{vlen\_le}, type-I automorphy, and type-II Kato input are
built with $\Psi=\texttt{Chr.Psi}(\texttt{Chr.dirG}\ \mathbb Q\ \psi)$.
Thus (GK.8), including its local-condition identification, is not
an instance of that theorem for $j=1,2$.
At an inert Frobenius a type-II block
$1\oplus(\chi^j\otimes W|_{G_K})$ has the same trace as
$1\oplus W|_{G_K}$.  This supplies the required trace comparison
\emph{if} the modular descent, level and other source admissibility
conditions for $W$ have been established.  It does not prove the
vanishing or integral comparison needed to extract that block
from the cubic Selmer group.

Finally, \path{HeightAndVanishing/Counts.lean}, record
\texttt{GD}, has twelve explicit conditions: $\ell\ne2$,
$e\ge1$, $N$ even, $\ell\nmid\operatorname{ord}\psi$,
$\operatorname{comp}_\ell(\psi)^2=1$,
$\psi(-1)=(-1)^e$, $[K:\mathbb Q]=2$, $K$ totally complex,
$\ell\nmid|D_K|h_K$, $(|D_K|,N)=1$, its stated noncontainment
of $\ker\overline\kappa_\ell\cap\ker(\texttt{chiP}_\ell\psi)$
in $G_K$, and $\ell-1\nmid2e$.
They are not omitted when using the source \texttt{PropGK}.
No replacement conditions for the cubic packet, no estimate for
their exceptional-prime contribution, and no cubic version of
(GK.8) are asserted proved here.  AC-K29 locates the remaining
arithmetic work in that extraction and those exceptions; the
common-integer size assembly adds no new growing-conductor
estimate once its exact hypotheses are available.
